\section*{الفصل \ref{ch:g12:subduction-magmatism} --- الصهارية في نطاقات الاندساس}

\begin{solution}{exo:g12:subduction-magmatism:1}
(أ) \omterm{def:g12:subduction-magmatism:arc}{قوس بركاني}؛ (ب) \omterm{def:g12:subduction-magmatism:hydrous-melting}{الانصهار المائي} (الانصهار بإضافة الماء)؛ (ج) \omterm{def:g12:subduction-magmatism:granodiorite}{الغرانوديوريت}؛
(د) \omterm{def:g12:subduction-magmatism:batholith}{الباثوليث}؛ (ه) \omterm{def:g12:subduction-magmatism:continental-growth}{نمو القارات}.
\end{solution}

\begin{solution}{exo:g12:subduction-magmatism:2}
(1) على الأقل $24/0.25 = \qty{96}{km/h}$، أي $\num{96000}/3600 \approx
\qty{27}{m/s}$. (2) كان الانفجار أسرع بما لا يقل عن $480/96 = 5$ مرات.
\end{solution}

\begin{solution}{exo:g12:subduction-magmatism:3}
(1) $838/1214 \approx 0.690$، أي نحو \qty{69}{\percent}. (2) على جانب
الصفيحة الراكبة، التي لا تغوص. (3) بين الخندق
والجبهة لا يكون \omterm{def:g11:subduction-zones:slab}{اللوح المندس} قد بلغ بعد العمق الكافي: فهو لم يصل إلى العمق،
نحو مئة كيلومتر، الذي يتحرر عنده ماؤه في وشاح ساخن
بما يكفي لينصهر ما إن يبتل.
\end{solution}

\begin{solution}{exo:g12:subduction-magmatism:4}
(1) $45/(45 + 15) = 0.75$: أي إن \qty{75}{\percent} من الفلسبار
بلاجيوكليز. (2) نعم: فهو خشن الحبيبات، ويحوي أكثر من
\qty{20}{\percent} من الكوارتز، وكسر البلاجيوكليز فيه يقع بين 65
و\qty{90}{\percent}.
\end{solution}

\begin{solution}{exo:g12:subduction-magmatism:5}
صهارات الأقواس أغنى بالسيليكا، فهي أكثر لزوجة؛ وتحمل
الماء الذي صنعها، فهي غنية بالغاز. واللزوجة العالية
ونسبة الغاز العالية كلتاهما تحابيان الثورانات الانفجارية.
\end{solution}

\begin{solution}{exo:g12:subduction-magmatism:6}
(1) $H = 150 \times \tan\ang{35} \approx 150 \times 0.700 \approx
\qty{105}{km}$. (2) إنه المتوسط العالمي.
\end{solution}

\begin{solution}{exo:g12:subduction-magmatism:7}
$d = H/\tan\theta$: $105/\tan\ang{20} \approx 105/0.364 \approx
\qty{288}{km}$ عند \ang{20}، و$105/\tan\ang{45} = \qty{105}{km}$ عند
\ang{45}. فكلما كان \omterm{def:g11:subduction-zones:slab}{اللوح المندس} أخف ميلًا، قام القوس أبعد عن الخندق.
\end{solution}

\begin{solution}{exo:g12:subduction-magmatism:8}
(1) $T = -5.104 \times 3.16^2 + 132.9 \times 3.16 + 1120.7 \approx -51.0 +
420.0 + 1120.7 \approx \qty{1490}{\celsius}$. (2) بنحو $1490 - 1000 =
\qty{490}{\celsius}$. (3) جافًا، لا: فالدرجة \qty{1300}{\celsius} أدنى من
\qty{1490}{\celsius}. ورطبًا، نعم: فهي فوق \qty{1000}{\celsius}.
\end{solution}

\begin{solution}{exo:g12:subduction-magmatism:9}
(1) $2.5/2.7 \approx 0.93$ من السنة، أي نحو 340 يومًا. (2) الانهيار نقل
صخرًا قديمًا كان موجودًا أصلًا؛ أما الأقواس فتصنع صخرًا جديدًا من الوشاح. فلا
يُقارن إلا الحجمان.
\end{solution}

\begin{solution}{exo:g12:subduction-magmatism:10}
(1) $2.7 \times \num{1000000} = \qty{2.7e6}{km^3}$. (2) $\num{9e9}/2.7
\approx \num{3.3e9}$ سنة، أي 3.3 مليار سنة. (3) هذا أقل من عمر
الأرض، فكان في وسع الأقواس أن تصنع \omterm{def:g12:continental-crust:continental-crust}{القشرة القارية} كلها؛ لكن
\omterm{def:g11:subduction-zones:subduction}{الاندساس} يعيد أيضًا جزءًا من القشرة، ففي الواقع صُنع كثير من القشرة
ودُمّر أكثر من مرة.
\end{solution}

\begin{solution}{exo:g12:subduction-magmatism:11}
(1) الزلازل الضحلة تقع أيضًا في الصفيحة الراكبة وعلى الصدع
الفاصل بين الصفيحتين؛ ولا تقع في \omterm{def:g11:subduction-zones:slab}{اللوح المندس} وحده إلا الزلازل العميقة.
(2) $116 - 80 = \qty{36}{km}$ على مسافة \qty{86.4}{km}: $\tan\theta = 36/86.4
\approx 0.417$، و$\theta \approx \ang{22.6}$.
\end{solution}

\begin{solution}{exo:g12:subduction-magmatism:12}
في المثلث القائم الوارد في القضية، \omterm{def:g11:subduction-zones:slab}{اللوح المندس} هو الوتر $L$
و$d$ الضلع المجاور للزاوية $\theta$، إذن $\cos\theta = d/L$ و$L =
d/\cos\theta$. ولليابان، $L = 270/\cos\ang{20} \approx 270/0.940 \approx
\qty{287}{km}$. وبمعدل \qty{83}{mm} في السنة: $\num{2.87e8}/83 \approx \num{3.5e6}$
سنة، أي نحو 3.5 مليون سنة.
\end{solution}

\begin{solution}{exo:g12:subduction-magmatism:13}
(1) $V = \pi r^2 h = \pi \times 0.5^2 \times 8 \approx \qty{6.3}{km^3}$.
(2) نحو $6.3/2.5 \approx 2.5$ ضعف حجم الانهيار.
(3) $P = \rho g z$: على عمق \qty{6}{km}، $2830 \times 9.81 \times 6000 \approx
\qty{1.67e8}{Pa} = \qty{0.17}{GPa}$؛ وعلى عمق \qty{14}{km}، نحو
\qty{3.89e8}{Pa} = \qty{0.39}{GPa}.
\end{solution}

\begin{solution}{exo:g12:subduction-magmatism:14}
(1) $2.7 - 3.2 = \qty{-0.5}{km^3}$ في السنة، أي خسارة \qty{5e5}{km^3} في كل
مليون سنة. (2) $500 \times \num{5e5} = \qty{2.5e8}{km^3}$، أي
$\num{2.5e8}/\num{9e9} \approx \qty{2.8}{\percent}$ من القشرة. (3) كل
معدل تقدير، وعدم اليقين فيه أكبر بكثير من الفرق البالغ
\qty{0.5}{km^3} في السنة بينهما: ففي حدود الأخطاء، يتوازن المعدلان.
\end{solution}

\begin{solution}{exo:g12:subduction-magmatism:15}
(1) $\num{20000} \times 10 = \qty{2e5}{km^3}$. (2) $\num{2e5}/1 =
\num{200000}$ سنة لكل الأقواس القارية في العالم مجتمعة. والقوس الواحد
لا يتلقى إلا جزءًا صغيرًا من تلك الصهارة، \omterm{def:g12:subduction-magmatism:batholith}{فالباثوليث} الواحد يستغرق
وقتًا أطول بكثير: فهو يُبنى \omterm{def:g11:igneous-rocks:intrusion}{بلوتونًا} بعد \omterm{def:g11:igneous-rocks:intrusion}{بلوتون}، على مدى ملايين السنين. (3) على
الأقل \qty{5}{km} من الصخر الذي فوق أعلاه.
\end{solution}

\begin{solution}{pb:g12:subduction-magmatism:1}
\textbf{1.} $111.2 \times \cos\ang{39} \approx 111.2 \times 0.777 \approx
\qty{86.4}{km}$.

\textbf{2.} $(144.1 - 141.0) \times 86.4 \approx 3.1 \times 86.4 \approx
\qty{268}{km}$، أي نحو \qty{270}{km}.

\textbf{3.} $(141.0 - 139.7) \times 86.4 \approx \qty{112}{km}$.

\textbf{4.} على الصفيحة الراكبة، التي لا تغوص، غرب
الخندق: فصفيحة المحيط الهادئ تغوص تحتها.

\textbf{5.} هناك، لا يكون \omterm{def:g11:subduction-zones:slab}{اللوح المندس} قد بلغ بعد العمق الكافي ليطلق ماءه في
وشاح ساخن بما يكفي لينصهر.

\textbf{6.} $116 - 80 = \qty{36}{km}$.

\textbf{7.} درجة طول واحدة، \qty{86.4}{km}.

\textbf{8.} $\tan\theta = 36/86.4 \approx 0.417$، إذن $\theta \approx
\ang{22.6}$.

\textbf{9.} يرتفع الخط \qty{36}{km} كل \qty{86.4}{km} نحو الشرق، فهو
يبلغ السطح على بعد $98/0.417 \approx \qty{235}{km}$ شرق إيواتي، عند
$141.0 + 235/86.4 \approx \ang{143.7}$ شرقًا: أي على بعد نحو $(144.1 - 143.7) \times 86.4
\approx \qty{35}{km}$ غرب الخندق.

\textbf{10.} \omterm{def:g11:subduction-zones:slab}{اللوح المندس} ليس مستويًا: فهو يدخل الخندق شبه
أفقي، ثم ينثني إلى الأسفل أكثر فأكثر كلما غاص.

\textbf{11.} $268 \times \tan\ang{20} \approx 268 \times 0.364 \approx
\qty{98}{km}$.

\textbf{12.} إنهما متفقتان: \qty{98}{km} بالطريقتين.

\textbf{13.} كلتاهما ضمن المدى العالمي، وأضحل قليلًا من
المتوسط \qty{105}{km}.

\textbf{14.} تقع ميغاتا على بعد \qty{112}{km} غرب إيواتي، فعمق \omterm{def:g11:subduction-zones:slab}{اللوح المندس} تحتها
نحو $98 + 0.417 \times 112 \approx \qty{145}{km}$.

\textbf{15.} في الوشاح، على نحو $98 - 35 = \qty{63}{km}$ تحت قاعدة
القشرة.

\textbf{16.} دخل قرب الظهرة التي تكوّنت عندها القشرة، من ماء
البحر؛ ويُحبس كثير منه في السربنتين.

\textbf{17.} $T = -5.104 \times 3.16^2 + 132.9 \times 3.16 + 1120.7 \approx
\qty{1490}{\celsius}$.

\textbf{18.} بمقدار يتراوح بين $1490 - 1000 = 490$ و$1490 - 950 = \qty{540}{\celsius}$:
أي نحو \qty{500}{\celsius}.

\textbf{19.} \omterm{def:g11:subduction-zones:slab}{اللوح المندس} أبرد بعدة مئات من الدرجات من الوشاح
المحيط به، وأبرد من أن ينصهر حتى مع الماء: فإسفين الوشاح الساخن الذي فوقه
هو الذي ينصهر ما إن يبتل.

\textbf{20.} يقع \omterm{def:g11:subduction-zones:slab}{اللوح المندس} على عمق نحو \textbf{\qty{100}{km} تحت الجبهة
البركانية} لشمال شرق اليابان (\qty{98}{km} بالطريقتين كلتيهما).
\end{solution}
